\documentclass{article}
\usepackage[utf8]{inputenc}
\usepackage[T1]{fontenc}
\usepackage{geometry}
\geometry{
    a4paper,
    left=2.5cm,
    right=2.5cm,
    top=2.5cm,
    bottom=2.5cm
}
\usepackage{graphicx}
\usepackage{float}
\usepackage{hyperref}
\usepackage[english]{babel}
\usepackage{amsmath}
\usepackage{amssymb}
\usepackage{array}
\usepackage{booktabs}
\usepackage{caption}
\usepackage{subcaption}
\usepackage{enumitem}
\usepackage{fancyhdr}
\usepackage{titlesec}
\usepackage{url}
\usepackage{textcomp}
\usepackage{xcolor}
\usepackage{listings}
\usepackage{multirow}
\usepackage{longtable}
\usepackage{tabularx}
\usepackage{times}

% Set up page style
\pagestyle{fancy}
\fancyhf{}
\rfoot{\thepage}
\renewcommand{\headrulewidth}{0pt}

% Title formatting
\titleformat{\chapter}[display]{\normalfont\bfseries}{\chaptertitle\thechapter}{0pt}{\Huge}
\titlespacing*{\chapter}{0pt}{50pt}{40pt}

% Document info
\title{Template Document}
\author{Academic Support Agent}
\date{\today}

\begin{document}
\maketitle

\preprint{APS/123-QED}

\title{T\'itulo de la pr\'actica}% Force line breaks with \\
%\thanks{A footnote to the article title}%

\author{Apellido1, Nombre1.}
\author{Apellido2, Nombre2.}
\author{Apellido3, Nombre3.}
 \affiliation{Universitdad Tecnol\'ogica de M\'exico, Campus Le\'on.}%Lines break automatically or can be forced with \\
 %t\email{Second.Author@myuniinstitution.edu}

\date{\today}% It is always \today, today,
             %  but any date may be explicitly specified

\begin{abstract}
Es un enunciado resumido del  contenido de  la  practica y  organizado secuencialmente  (menos de 100 palabras). Establece brevemente el problema y prop sito de la pr ctica. Indica el plan te rico o experimental a que se sigui . Resume las principales aportaciones y principales conclusiones.\\ \textbf{Importante}:El  resumen  es  la  informaci n  condensada  que  el  profesor  recibe  de  la  pr ctica realizada.  Debe  ser  corto  y  claro  para  que  el  profesor  determine  si  se  entendi   y  realizo correctamente  la  pr ctica.  Incluye  informaci n  de  seguridad  cuando  sea  necesario.  No  eval es conclusiones.\\\textbf{S per  Importante}:Escribe  el  resumen  al   ltimo  para  estar  seguro  que  refleja  con  exactitud  la pr ctica realizada
\begin{description}
\item[Palabras clave]: medici n, vernier.
\end{description}
\end{abstract}

%\keywords{Suggested keywords}%Use showkeys class option if keyword
                              %display desired
\maketitle

%

\section{\label{sec:introduccion}Introducci n}

 La  introducci n  debe  responder  a  la  pregunta  de  "porqu   se  ha  hecho  este  trabajo".Una  buena introducci n,  es  una  oraci n  clara  del  problema  y  de  las  razones  por  las  que  lo  estamos estudiando.  Nos  da  una  concisa  y  apropiada  discusi n  del  problema,  su  significado,  alcances  y limitaciones.\\\textbf{Importante}:En esta secci n se responde a la pregunta de "c mo se ha hecho la pr ctica".\\\textbf{S per Importante}:Es conveniente que el  ltimo p rrafo de la Introducci n se utilice para resumir el objetivode la pr ctica.

\section{Objetivo y planteamiento del problema}

Identificaci n del problema central de estudio y enfoque de la pr ctica.Importante:Esta  secci n  se  relaciona  con  el  procedimiento  experimental  presentado  en  el protocolo, pero no lo es del todo. Se trata que identifiques dentro de tu desarrollo experimental el sentido que tuvo la pr ctica para poder abordar el tema a estudiar.

\section{Desarrollo experimental}

En esta secci\'on deben describirse los materiales y m todos utilizados para el desarrollo de la pr ctica. Para describir el arreglo experimental, puede utilizarse una lista:
\begin{itemize}
    \item Elemento 1
    \item Elemento 2
    \item \ldots
\end{itemize}
o puede describirse en un p rrafo c mo se form  el arreglo. Sin embargo, es de car cter obligatorio que el arreglo experimental se describa en un diagrama (no fotos).
Finalmente, en esta secci\'on debe describirse c mo se llev  a cabo el experimento.
\section{Discusi n y Resultados}
Conjunto  de  datos  obtenidos  experimentalmente  y  tratados  estad sticamente.  Usa  tablas  paraorganizar y resumir los resultados\\\textbf{Importante}:Incluye  solo  los  datos  importantes  y  relevantes,  pero  suficientes  para  justificar  tus conclusiones. Usa ecuaciones, gr ficos y figuras.\\\textbf{S per  Importante}:Si  cuentas  con  datos  te ricos  no  olvides  obtener  el  error  experimental  (\%E). Esta informaci n es muy relevante para justificar tus conclusiones.
Cuando construyas tablas no olvides:
\begin{enumerate}
    \item  Titulo de la tabla
    \item Distinci n clara de celdas.
    \item T tulos de columnas.
    \item Notas al pie de la tabla
    \item Numeraci n de tablas
\end{enumerate}
\section{Conclusiones}
Las conclusiones en un reporte deber n recoger los puntos primordiales del trabajo realizado, siendo una de las partes de un informe importante.

\section{Bibliograf a}

En esta secci n se anotan todas las fuentes bibliogr ficas consultadas para el desarrollo de la pr ctica.

%% Produces the bibliography via BibTeX.
\end{document}